% GENERATED FILE. Edit canonical lessons, then run npm run generate:books.
% canonical-source-hash: fnv1a64:f86c8ad27862e0e6
% canonical-lessons: MW-C283-moti, MW-C283-bidhno, MW-C283-jarno, MW-C283-kasido, MW-C283-sajavno

\chapter{Pearls and Embroidery}
\label{ch:mw-pearls-and-embroidery}
\hlchaptermodality{283}

\begin{chapteropening}
\textbf{By the end of this chapter:} \emph{I can name crafts I like: beads, inlay and embroidery.}

Complete the last lesson of chapter 283: I can name crafts I like: beads, inlay and embroidery.
\end{chapteropening}

\section[\texorpdfstring{\mw{मोती} (motī)}{मोती (motī)}]{\mw{मोती} (motī) — a pearl}

\label{lesson:MW-C283-moti}



\begin{quote}
Before the new one: say the Marwadi for a garland, then the Marwadi for to string beads.
\end{quote}

\subsection*{You'll want to know: \mw{मोती}}
\textbf{\mw{मोती}} — \emph{motī} — \textquotedblleft{}a pearl\textquotedblright{}.

\begin{quote}
\mw{यो} \mw{मोती} \mw{घणो} \mw{महंगो} \mw{है।}
\end{quote}

\begin{quote}
\emph{yo motī ghaṇo mahaṅgo hai.}
\end{quote}

\begin{quote}
This pearl is very expensive.
\end{quote}

A masculine noun.

\begin{cousinweb}
\textbf{\mw{मोती}} (\emph{motī}) is inherited from Sanskrit \emph{mauktika}, \textquotedblleft{}a pearl\textquotedblright{}; CDIAL lists Old Marwari \emph{motī}, \textquotedblleft{}pearl\textquotedblright{} (CDIAL 10365).
\end{cousinweb}

\begin{grammarlens}[title={What you've built}]
One more word for this chapter.
\end{grammarlens}

\subsection*{Your turn}
\begin{itemize}
\raggedright
  \item \emph{Say it:} \emph{motī}
  \item \emph{Say it:} \emph{motī}, once more
  \item \emph{Say it:} \emph{yo motī ghaṇo mahaṅgo hai.} — This pearl is very expensive.
  \item \emph{Recall:} say the Marwadi for a garland, then the Marwadi for to string beads, then say \emph{motī} again
\end{itemize}

\begin{tcolorbox}[breakable,colback=teal!4,colframe=teal!35!black,title={Before you move on}]
What does \mw{मोती} mean? (\textquotedblleft{}A pearl\textquotedblright{}.) Say it once more.
\end{tcolorbox}
\section[\texorpdfstring{\mw{बींधणो} (bī̃dhṇo)}{बींधणो (bī̃dhṇo)}]{\mw{बींधणो} (bī̃dhṇo) — to pierce (a bead)}

\label{lesson:MW-C283-bidhno}



\begin{quote}
Before the new one: say the Marwadi for to string beads, then the Marwadi for a pearl.
\end{quote}

\subsection*{You'll want to know: \mw{बींधणो}}
\textbf{\mw{बींधणो}} — \emph{bī̃dhṇo} — \textquotedblleft{}to pierce (a bead)\textquotedblright{}.

\begin{quote}
\mw{मां} \mw{मोती} \mw{बींधै} \mw{है।}
\end{quote}

\begin{quote}
\emph{mā̃ motī bī̃dhai hai.}
\end{quote}

\begin{quote}
Mother pierces the pearls.
\end{quote}

Past (I …): \mw{म्हैं} \mw{मोती} \mw{बींध्यो}, \emph{mhaĩ motī bī̃dhyo}.

\begin{cousinweb}
\textbf{\mw{बींधणो}} (\emph{bī̃dhṇo}) is inherited from Sanskrit \emph{viddha}, \textquotedblleft{}pierced\textquotedblright{}; CDIAL lists Old Marwari \emph{abīdha}, \textquotedblleft{}unpierced\textquotedblright{}, and Hindi \emph{bīdhnā}, \textquotedblleft{}to pierce\textquotedblright{} (CDIAL 11739).
\end{cousinweb}

\begin{grammarlens}[title={What you've built}]
One more word for this chapter.
\end{grammarlens}

\subsection*{Your turn}
\begin{itemize}
\raggedright
  \item \emph{Say it:} \emph{bī̃dhṇo}
  \item \emph{Say it:} \emph{bī̃dhṇo}, once more
  \item \emph{Say it:} \emph{mā̃ motī bī̃dhai hai.} — Mother pierces the pearls.
  \item \emph{Recall:} say the Marwadi for to string beads, then the Marwadi for a pearl, then say \emph{bī̃dhṇo} again
\end{itemize}

\begin{tcolorbox}[breakable,colback=teal!4,colframe=teal!35!black,title={Before you move on}]
What does \mw{बींधणो} mean? (\textquotedblleft{}To pierce (a bead)\textquotedblright{}.) Say it once more.
\end{tcolorbox}
\section[\texorpdfstring{\mw{जड़णो} (jaṛṇo)}{जड़णो (jaṛṇo)}]{\mw{जड़णो} (jaṛṇo) — to set (a stone); to inlay}

\label{lesson:MW-C283-jarno}



\begin{quote}
Before the new one: say the Marwadi for a pearl, then the Marwadi for to pierce.
\end{quote}

\subsection*{You'll want to know: \mw{जड़णो}}
\textbf{\mw{जड़णो}} — \emph{jaṛṇo} — \textquotedblleft{}to set (a stone); to inlay\textquotedblright{}.

\begin{quote}
\mw{काको} \mw{अंगूठी} \mw{माथै} \mw{मोती} \mw{जड़ै} \mw{है।}
\end{quote}

\begin{quote}
\emph{kāko aṅgūṭhī māthai motī jaṛai hai.}
\end{quote}

\begin{quote}
My uncle sets a pearl on a ring.
\end{quote}

Past (I …): \mw{म्हैं} \mw{मोती} \mw{जड़्यो}, \emph{mhaĩ motī jaṛyo}.

\begin{cousinweb}
\textbf{\mw{जड़णो}} (\emph{jaṛṇo}) is inherited from a reconstructed \emph{*jaḍati}, \textquotedblleft{}joins, sets\textquotedblright{}, which CDIAL takes from Dravidian; it lists Old Marwari \emph{jaṛāū}, \textquotedblleft{}inlaid\textquotedblright{} (CDIAL 5091).
\end{cousinweb}

\begin{grammarlens}[title={What you've built}]
One more word for this chapter.
\end{grammarlens}

\subsection*{Your turn}
\begin{itemize}
\raggedright
  \item \emph{Say it:} \emph{jaṛṇo}
  \item \emph{Say it:} \emph{jaṛṇo}, once more
  \item \emph{Say it:} \emph{kāko aṅgūṭhī māthai motī jaṛai hai.} — My uncle sets a pearl on a ring.
  \item \emph{Recall:} say the Marwadi for a pearl, then the Marwadi for to pierce, then say \emph{jaṛṇo} again
\end{itemize}

\begin{tcolorbox}[breakable,colback=teal!4,colframe=teal!35!black,title={Before you move on}]
What does \mw{जड़णो} mean? (\textquotedblleft{}To set (a stone); to inlay\textquotedblright{}.) Say it once more.
\end{tcolorbox}
\section[\texorpdfstring{\mw{कसीदो} (kasīdo)}{कसीदो (kasīdo)}]{\mw{कसीदो} (kasīdo) — embroidery}

\label{lesson:MW-C283-kasido}



\begin{quote}
Before the new one: say the Marwadi for to pierce, then the Marwadi for to inlay.
\end{quote}

\subsection*{You'll want to know: \mw{कसीदो}}
\textbf{\mw{कसीदो}} — \emph{kasīdo} — \textquotedblleft{}embroidery\textquotedblright{}.

\begin{quote}
\mw{यो} \mw{कसीदो} \mw{सुंदर} \mw{है।}
\end{quote}

\begin{quote}
\emph{yo kasīdo sundar hai.}
\end{quote}

\begin{quote}
This embroidery is beautiful.
\end{quote}

A masculine noun.

\begin{cousinweb}
\textbf{\mw{कसीदो}} (\emph{kasīdo}) is usually traced to Persian \emph{kashīda}, \textquotedblleft{}drawn out\textquotedblright{}, used of embroidered work; Marwadi gives it the -\mw{ो} ending of a masculine noun.
\end{cousinweb}

\begin{grammarlens}[title={What you've built}]
One more word for this chapter.
\end{grammarlens}

\subsection*{Your turn}
\begin{itemize}
\raggedright
  \item \emph{Say it:} \emph{kasīdo}
  \item \emph{Say it:} \emph{kasīdo}, once more
  \item \emph{Say it:} \emph{yo kasīdo sundar hai.} — This embroidery is beautiful.
  \item \emph{Recall:} say the Marwadi for to pierce, then the Marwadi for to inlay, then say \emph{kasīdo} again
\end{itemize}

\begin{tcolorbox}[breakable,colback=teal!4,colframe=teal!35!black,title={Before you move on}]
What does \mw{कसीदो} mean? (\textquotedblleft{}Embroidery\textquotedblright{}.) Say it once more.
\end{tcolorbox}
\section[\texorpdfstring{\mw{सजावणो} (sajāvṇo)}{सजावणो (sajāvṇo)}]{\mw{सजावणो} (sajāvṇo) — to decorate}

\label{lesson:MW-C283-sajavno}



\begin{quote}
Before the new one: say the Marwadi for to inlay, then the Marwadi for embroidery.
\end{quote}

\subsection*{You'll want to know: \mw{सजावणो}}
\textbf{\mw{सजावणो}} — \emph{sajāvṇo} — \textquotedblleft{}to decorate\textquotedblright{}.

\begin{quote}
\mw{छोरी} \mw{घर} \mw{सजावै} \mw{है।}
\end{quote}

\begin{quote}
\emph{chhorī ghar sajāvai hai.}
\end{quote}

\begin{quote}
The girl decorates the house.
\end{quote}

Past (I …): \mw{म्हैं} \mw{घर} \mw{सजायो}, \emph{mhaĩ ghar sajāyo}.

\begin{cousinweb}
\textbf{\mw{सजावणो}} (\emph{sajāvṇo}) is inherited from Sanskrit \emph{sajjayati}, \textquotedblleft{}equips, prepares\textquotedblright{}; CDIAL lists Old Marwari \emph{sājaï}, \textquotedblleft{}adorns\textquotedblright{}, and Hindi \emph{sajānā}, \textquotedblleft{}to put in order, decorate\textquotedblright{} (CDIAL 13091).
\end{cousinweb}

\begin{grammarlens}[title={What you've built}]
One more word for this chapter.
\end{grammarlens}

\subsection*{Your turn}
\begin{itemize}
\raggedright
  \item \emph{Say it:} \emph{sajāvṇo}
  \item \emph{Say it:} \emph{sajāvṇo}, once more
  \item \emph{Say it:} \emph{chhorī ghar sajāvai hai.} — The girl decorates the house.
  \item \emph{Recall:} say the Marwadi for to inlay, then the Marwadi for embroidery, then say \emph{sajāvṇo} again
\end{itemize}

\begin{tcolorbox}[breakable,colback=teal!4,colframe=teal!35!black,title={Before you move on}]
What does \mw{सजावणो} mean? (\textquotedblleft{}To decorate\textquotedblright{}.) Say it once more.
\end{tcolorbox}
